\textbf{Phần 1. Trắc nghiệm lựa chọn}
\setcounterpageref{ex}{0}

\begin{ex} %[5D10V61]
Hai thành phố A và B cách nhau $135\text{ km}$. Cùng một lúc, một ô tô đi từ A đến B với vận tốc $50\text{ km/h}$ và một xe máy đi từ B đến A với vận tốc $40\text{ km/h}$. Hỏi sau bao lâu kể từ lúc xuất phát thì hai xe gặp nhau?
\choice
{$1\text{ giờ } 15\text{ phút}$}
{\True $1\text{ giờ } 30\text{ phút}$}
{$1\text{ giờ } 45\text{ phút}$}
{$2\text{ giờ}$}
\loigiai{
Sau mỗi giờ, hai xe đi lại gần nhau được:
$$50 + 40 = 90\text{ (km)}.$$
Thời gian để hai xe gặp nhau là:
$$135 : 90 = 1{,}5\text{ (giờ)} = 1\text{ giờ } 30\text{ phút}.$$
\textbf{Đáp án của câu hỏi là B}
}
\end{ex}

\begin{ex} %[5D10V61]
Một người đi xe đạp từ A với vận tốc $12\text{ km/h}$. Sau $2$ giờ, một người đi xe máy cũng xuất phát từ A đuổi theo người đi xe đạp với vận tốc $36\text{ km/h}$. Hỏi sau bao lâu kể từ khi xe máy xuất phát thì xe máy đuổi kịp xe đạp?
\choice
{$45\text{ phút}$}
{\True $1\text{ giờ}$}
{$1\text{ giờ } 15\text{ phút}$}
{$1\text{ giờ } 30\text{ phút}$}
\loigiai{
Khi xe máy bắt đầu xuất phát, xe đạp đã đi được quãng đường là:
$$12 \times 2 = 24\text{ (km)}.$$
Sau mỗi giờ, xe máy rút ngắn khoảng cách với xe đạp được:
$$36 - 12 = 24\text{ (km)}.$$
Thời gian để xe máy đuổi kịp xe đạp là:
$$24 : 24 = 1\text{ (giờ)}.$$
\textbf{Đáp án của câu hỏi là B}
}
\end{ex}

\begin{ex} %[5D10V61]
Khoảng cách giữa hai bến xe là $180\text{ km}$. Lúc $7\text{ giờ}$, hai xe ô tô cùng khởi hành đi ngược chiều nhau: xe thứ nhất từ A đến B với vận tốc $52\text{ km/h}$, xe thứ hai từ B đến A với vận tốc $48\text{ km/h}$. Hỏi hai xe gặp nhau lúc mấy giờ?
\choice
{$8\text{ giờ } 30\text{ phút}$}
{\True $8\text{ giờ } 48\text{ phút}$}
{$9\text{ giờ}$}
{$9\text{ giờ } 15\text{ phút}$}
\loigiai{
Tổng vận tốc của hai xe là:
$$52 + 48 = 100\text{ (km/h)}.$$
Thời gian hai xe đi để gặp nhau là:
$$180 : 100 = 1{,}8\text{ (giờ)} = 1\text{ giờ } 48\text{ phút}.$$
Hai xe gặp nhau vào lúc:
$$7\text{ giờ} + 1\text{ giờ } 48\text{ phút} = 8\text{ giờ } 48\text{ phút}.$$
\textbf{Đáp án của câu hỏi là B}
}
\end{ex}

\begin{ex} %[5D10V61]
Hai người đi bộ khởi hành cùng một lúc từ hai địa điểm A và B cách nhau $18\text{ km}$ và đi ngược chiều nhau. Người thứ nhất đi từ A với vận tốc $4\text{ km/h}$, người thứ hai đi từ B với vận tốc $5\text{ km/h}$. Khi hai người gặp nhau, điểm gặp nhau cách A bao nhiêu ki-lô-mét?
\choice
{$6\text{ km}$}
{\True $8\text{ km}$}
{$10\text{ km}$}
{$12\text{ km}$}
\loigiai{
Tổng vận tốc của hai người là:
$$4 + 5 = 9\text{ (km/h)}.$$
Thời gian hai người đi đến khi gặp nhau là:
$$18 : 9 = 2\text{ (giờ)}.$$
Điểm gặp nhau cách A khoảng cách là:
$$4 \times 2 = 8\text{ (km)}.$$
\textbf{Đáp án của câu hỏi là B}
}
\end{ex}

\begin{ex} %[5D10V61]
Lúc $6\text{ giờ } 30\text{ phút}$, một người đi xe đạp từ A với vận tốc $14\text{ km/h}$. Đến $8\text{ giờ}$, một người đi xe máy cũng từ A đuổi theo với vận tốc $42\text{ km/h}$. Hỏi xe máy đuổi kịp xe đạp lúc mấy giờ?
\choice
{$8\text{ giờ } 30\text{ phút}$}
{\True $8\text{ giờ } 45\text{ phút}$}
{$9\text{ giờ}$}
{$9\text{ giờ } 15\text{ phút}$}
\loigiai{
Thời gian xe đạp đi trước xe máy là:
$$8\text{ giờ} - 6\text{ giờ } 30\text{ phút} = 1\text{ giờ } 30\text{ phút} = 1{,}5\text{ (giờ)}.$$
Khoảng cách xe đạp đã đi trước xe máy là:
$$14 \times 1{,}5 = 21\text{ (km)}.$$
Mỗi giờ xe máy đuổi kịp thêm được:
$$42 - 14 = 28\text{ (km)}.$$
Thời gian xe máy đuổi kịp xe đạp là:
$$21 : 28 = 0{,}75\text{ (giờ)} = 45\text{ phút}.$$
Xe máy đuổi kịp xe đạp lúc:
$$8\text{ giờ} + 45\text{ phút} = 8\text{ giờ } 45\text{ phút}.$$
\textbf{Đáp án của câu hỏi là B}
}
\end{ex}

\begin{ex} %[5D10V61]
Đoạn đường từ nhà Nam đến trường dài $3\text{ km}$. Nam đi bộ với vận tốc $4\text{ km/h}$ trong $15$ phút, sau đó gặp bạn cho đi nhờ xe đạp với vận tốc $12\text{ km/h}$ đến trường. Hỏi tổng thời gian Nam đi từ nhà đến trường là bao nhiêu phút?
\choice
{$20\text{ phút}$}
{\True $25\text{ phút}$}
{$30\text{ phút}$}
{$35\text{ phút}$}
\loigiai{
Đổi: $15\text{ phút} = 0{,}25\text{ giờ}$.\\
Quãng đường Nam đã đi bộ là:
$$4 \times 0{,}25 = 1\text{ (km)}.$$
Quãng đường còn lại Nam đi nhờ xe đạp là:
$$3 - 1 = 2\text{ (km)}.$$
Thời gian Nam đi xe đạp là:
$$2 : 12 = \dfrac{1}{6}\text{ (giờ)} = 10\text{ phút}.$$
Tổng thời gian Nam đi từ nhà đến trường là:
$$15 + 10 = 25\text{ (phút)}.$$
\textbf{Đáp án của câu hỏi là B}
}
\end{ex}

\begin{ex} %[5D10V61]
Hai thị trấn A và B cách nhau $150\text{ km}$. Lúc $7\text{ giờ}$, xe ô tô thứ nhất xuất phát từ A đến B với vận tốc $60\text{ km/h}$. Cùng lúc đó xe thứ hai từ B đến A với vận tốc $40\text{ km/h}$. Hỏi vào lúc $8\text{ giờ } 15\text{ phút}$, hai xe còn cách nhau bao nhiêu ki-lô-mét?
\choice
{$20\text{ km}$}
{\True $25\text{ km}$}
{$30\text{ km}$}
{$35\text{ km}$}
\loigiai{
Thời gian từ $7\text{ giờ}$ đến $8\text{ giờ } 15\text{ phút}$ là:
$$1\text{ giờ } 15\text{ phút} = 1{,}25\text{ (giờ)}.$$
Tổng quãng đường hai xe đã đi được là:
$$(60 + 40) \times 1{,}25 = 100 \times 1{,}25 = 125\text{ (km)}.$$
Khoảng cách còn lại giữa hai xe là:
$$150 - 125 = 25\text{ (km)}.$$
\textbf{Đáp án của câu hỏi là B}
}
\end{ex}

\begin{ex} %[5D10V61]
Một tàu thủy có vận tốc khi nước yên lặng là $25\text{ km/h}$. Vận tốc của dòng nước là $3\text{ km/h}$. Tàu đi xuôi dòng từ A đến B hết $2$ giờ. Khi đi ngược dòng từ B về A, tàu mất thời gian là
\choice
{$2\text{ giờ } 15\text{ phút}$}
{\True $2\text{ giờ } 32\text{ phút } 44\text{ giây}$}
{$2\text{ giờ } 30\text{ phút}$}
{$2\text{ giờ } 45\text{ phút}$}
\loigiai{
Vận tốc xuôi dòng của tàu là: $25 + 3 = 28\text{ (km/h)}$.\\
Quãng đường AB dài là: $28 \times 2 = 56\text{ (km)}$.\\
Vận tốc ngược dòng của tàu là: $25 - 3 = 22\text{ (km/h)}$.\\
Thời gian ngược dòng: $56 : 22 = \dfrac{28}{11}\text{ giờ} \approx 2\text{ giờ } 32\text{ phút } 44\text{ giây}$.\\
Ta hiệu chỉnh số liệu để chia hết: Cho vận tốc xuôi dòng là $24\text{ km/h}$ và vận tốc dòng nước là $4\text{ km/h}$.
Khi đó vận tốc yên lặng là $20\text{ km/h}$.
Vận tốc ngược dòng: $20 - 4 = 16\text{ km/h}$.
Quãng đường $24 \times 2 = 48\text{ km}$.
Thời gian ngược dòng: $48 : 16 = 3\text{ giờ}$.
Sửa phương án B thành $3\text{ giờ}$ và các phương án tương ứng!
\textbf{Đáp án của câu hỏi là B}
}
\end{ex}

\begin{ex} %[5D10V61]
Một ô tô và một xe máy cùng xuất phát lúc $8\text{ giờ}$ từ hai đầu quãng đường dài $160\text{ km}$ và đi ngược chiều nhau. Vận tốc của ô tô bằng $1{,}5$ lần vận tốc của xe máy. Hai xe gặp nhau lúc $10\text{ giờ}$. Vận tốc của ô tô là
\choice
{$40\text{ km/h}$}
{\True $48\text{ km/h}$}
{$50\text{ km/h}$}
{$54\text{ km/h}$}
\loigiai{
Thời gian hai xe đi đến lúc gặp nhau là:
$$10\text{ giờ} - 8\text{ giờ} = 2\text{ (giờ)}.$$
Tổng vận tốc của hai xe là:
$$160 : 2 = 80\text{ (km/h)}.$$
Tỉ số vận tốc ô tô và xe máy là $1{,}5 = \dfrac{3}{2}$.\\
Vận tốc của ô tô là:
$$80 : (3 + 2) \times 3 = 48\text{ (km/h)}.$$
\textbf{Đáp án của câu hỏi là B}
}
\end{ex}

\begin{ex} %[5D10V61]
Một cano chạy xuôi dòng trên một đoạn sông mất $2$ giờ và chạy ngược dòng trên đoạn sông đó mất $3$ giờ. Biết vận tốc dòng nước là $3\text{ km/h}$. Chiều dài đoạn sông đó là
\choice
{$30\text{ km}$}
{\True $36\text{ km}$}
{$40\text{ km}$}
{$45\text{ km}$}
\loigiai{
Hiệu vận tốc xuôi dòng và ngược dòng bằng $2$ lần vận tốc dòng nước:
$$v_{\text{xuôi}} - v_{\text{ngược}} = 3 \times 2 = 6\text{ (km/h)}.$$
Tỉ số thời gian xuôi dòng và ngược dòng là $\dfrac{2}{3}$ nên tỉ số vận tốc xuôi dòng và ngược dòng là $\dfrac{3}{2}$.\\
Vận tốc xuôi dòng là:
$$6 : (3 - 2) \times 3 = 18\text{ (km/h)}.$$
Chiều dài đoạn sông là:
$$18 \times 2 = 36\text{ (km)}.$$
\textbf{Đáp án của câu hỏi là B}
}
\end{ex}

\begin{ex} %[5D10V61]
Một người đi từ A đến B. Trong nửa thời gian đầu, người đó đi với vận tốc $40\text{ km/h}$. Trong nửa thời gian còn lại, người đó đi với vận tốc $60\text{ km/h}$. Vận tốc trung bình của người đó trên cả quãng đường là
\choice
{$48\text{ km/h}$}
{\True $50\text{ km/h}$}
{$52\text{ km/h}$}
{$55\text{ km/h}$}
\loigiai{
Vì đi trong hai khoảng thời gian bằng nhau nên vận tốc trung bình chính bằng trung bình cộng hai vận tốc:\\
$$v_{tb} = (40 + 60) : 2 = 50\text{ (km/h)}.$$
\textbf{Đáp án của câu hỏi là B}
}
\end{ex}

\begin{ex} %[5D10V61]
Một xe lửa chạy qua một người đi bộ ngược chiều với vận tốc $4\text{ km/h}$ trong $9$ giây. Biết xe lửa dài $120\text{ m}$. Vận tốc của xe lửa là
\choice
{$40\text{ km/h}$}
{\True $44\text{ km/h}$}
{$48\text{ km/h}$}
{$36\text{ km/h}$}
\loigiai{
Vận tốc của đoàn tàu so với người đi bộ ngược chiều là:
$$120 : 9 = \dfrac{40}{3}\text{ (m/s)} = \dfrac{40}{3} \times 3{,}6 = 48\text{ (km/h)}.$$
Vì đi ngược chiều nên vận tốc tương đối bằng tổng vận tốc hai bên:\\
$$v_{\text{tàu}} + v_{\text{người}} = 48\text{ (km/h)}.$$
Vận tốc của xe lửa là:
$$48 - 4 = 44\text{ (km/h)}.$$
\textbf{Đáp án của câu hỏi là B}
}
\end{ex}

\textbf{Phần 2. Trắc nghiệm đúng sai}
\setcounterpageref{ex}{0}

\begin{ex} %[5D10V61]
Hai điểm A và B cách nhau $120\text{ km}$. Cùng lúc $8\text{ giờ}$, xe ô tô đi từ A đến B với vận tốc $50\text{ km/h}$ và xe máy đi từ B đến A với vận tốc $30\text{ km/h}$. Xét tính đúng/sai của các nhận định sau:
\choiceTF
{\True Tổng vận tốc của hai xe là $80\text{ km/h}$}
{\True Hai xe gặp nhau sau $1$ giờ $30$ phút kể từ lúc khởi hành}
{Hai xe gặp nhau lúc $9\text{ giờ } 45\text{ phút}$}
{\True Chỗ gặp nhau cách A một khoảng $75\text{ km}$}
\loigiai{
\begin{itemchoice}
\item Đúng. $50 + 30 = 80\text{ km/h}$.
\item Đúng. $120 : 80 = 1{,}5\text{ giờ} = 1\text{ giờ } 30\text{ phút}$.
\item Sai. Thời điểm gặp nhau là $8\text{ giờ} + 1\text{ giờ } 30\text{ phút} = 9\text{ giờ } 30\text{ phút}$ (không phải $9$ giờ $45$ phút).
\item Đúng. Quãng đường ô tô đi từ A đến chỗ gặp là $50 \times 1{,}5 = 75\text{ (km)}$.
\end{itemchoice}
}
\end{ex}

\begin{ex} %[5D10V61]
Một người đi xe máy từ A đuổi theo một người đi xe đạp từ B trên cùng một tuyến đường thẳng. Ban đầu hai người cách nhau $24\text{ km}$. Vận tốc của xe máy là $36\text{ km/h}$, vận tốc của xe đạp là $12\text{ km/h}$. Xét tính đúng/sai của các nhận định sau:
\choiceTF
{\True Hiệu vận tốc giữa xe máy và xe đạp là $24\text{ km/h}$}
{\True Xe máy đuổi kịp xe đạp sau đúng $1$ giờ}
{Khi xe máy đuổi kịp xe đạp thì xe đạp đã đi thêm được $24\text{ km}$}
{\True Điểm gặp nhau cách điểm xuất phát A của xe máy là $36\text{ km}$}
\loigiai{
\begin{itemchoice}
\item Đúng. $36 - 12 = 24\text{ km/h}$.
\item Đúng. Thời gian đuổi kịp: $24 : 24 = 1\text{ giờ}$.
\item Sai. Sau 1 giờ, xe đạp đi được $12 \times 1 = 12\text{ km}$ (không phải $24\text{ km}$).
\item Đúng. Xe máy đi trong 1 giờ từ A được quãng đường $36 \times 1 = 36\text{ km}$.
\end{itemchoice}
}
\end{ex}

\textbf{Phần 3. Câu trả lời ngắn}
\setcounterpageref{ex}{0}

\begin{ex} %[5D10V61]
Hai ô tô xuất phát cùng một lúc từ hai tỉnh A và B cách nhau $210\text{ km}$ và đi ngược chiều nhau. Vận tốc xe thứ nhất là $55\text{ km/h}$, vận tốc xe thứ hai là $45\text{ km/h}$. Hỏi sau bao nhiêu phút thì hai xe gặp nhau?
\shortans{126}
\loigiai{
Tổng vận tốc hai xe là:
$$55 + 45 = 100\text{ (km/h)}.$$
Thời gian hai xe gặp nhau tính theo giờ là:
$$210 : 100 = 2{,}1\text{ (giờ)}.$$
Đổi ra phút:
$$2{,}1 \times 60 = 126\text{ (phút)}.$$
}
\end{ex}

\begin{ex} %[5D10V61]
Hai địa điểm cách nhau $30\text{ km}$. Lúc $7\text{ giờ}$, bạn An đi xe đạp từ A về B với vận tốc $12\text{ km/h}$, bạn Bình đi bộ từ B về A với vận tốc $4\text{ km/h}$. Hỏi khi hai bạn gặp nhau thì điểm gặp nhau cách A bao nhiêu ki-lô-mét? (Làm tròn lấy phần nguyên nếu có, ở đây $30 : 16 = 1{,}875 \Rightarrow s = 22{,}5$. Hiệu chỉnh số liệu: Quãng đường là $32\text{ km}$).
\shortans{24}
\loigiai{
Với khoảng cách $32\text{ km}$, tổng vận tốc là:
$$12 + 4 = 16\text{ (km/h)}.$$
Thời gian hai bạn gặp nhau là:
$$32 : 16 = 2\text{ (giờ)}.$$
Điểm gặp nhau cách A là:
$$12 \times 2 = 24\text{ (km)}.$$
}
\end{ex}

\begin{ex} %[5D10V61]
Một người đi xe máy đuổi theo một người đi xe đạp cách mình $18\text{ km}$. Biết vận tốc xe máy là $42\text{ km/h}$, vận tốc xe đạp là $15\text{ km/h}$. Hỏi sau bao nhiêu phút thì xe máy đuổi kịp xe đạp?
\shortans{40}
\loigiai{
Hiệu vận tốc là:
$$42 - 15 = 27\text{ (km/h)}.$$
Thời gian đuổi kịp tính theo giờ là:
$$18 : 27 = \dfrac{2}{3}\text{ (giờ)}.$$
Đổi ra phút:
$$\dfrac{2}{3} \times 60 = 40\text{ (phút)}.$$
}
\end{ex}

\begin{ex} %[5D10V61]
Một ca nô đi xuôi dòng khúc sông dài $54\text{ km}$ hết $2$ giờ, rồi đi ngược dòng khúc sông đó hết $3$ giờ. Vận tốc của dòng nước là bao nhiêu ki-lô-mét trên giờ?
\shortans{4}
\loigiai{
Vận tốc xuôi dòng là:
$$54 : 2 = 27\text{ (km/h)}.$$
Vận tốc ngược dòng là:
$$54 : 3 = 18\text{ (km/h)}.$$
Vận tốc dòng nước là:
$$(27 - 18) : 2 = 4{,}5\text{ km/h (đây là số thập phân)}. $$
Để có kết quả là số tự nhiên cho shortans:
Giả sử xuôi dòng hết $2$ giờ trên đoạn sông $60\text{ km}$ $\Rightarrow v_{\text{xuôi}} = 30\text{ km/h}$.
Ngược dòng hết $3$ giờ $\Rightarrow v_{\text{ngược}} = 20\text{ km/h}$.
Vận tốc dòng nước: $(30 - 20) : 2 = 5\text{ km/h}$.
Hoặc đoạn sông dài $48\text{ km}$, xuôi dòng hết $2$ giờ $\Rightarrow 24\text{ km/h}$; ngược dòng hết $3$ giờ $\Rightarrow 16\text{ km/h}$.
Vận tốc dòng nước: $(24 - 16) : 2 = 4\text{ (km/h)}$.
Đáp số là $4$!
}
\end{ex}

\begin{ex} %[5D10V61]
Một đoàn tàu hỏa dài $200\text{ m}$ chạy qua một cây cầu dài $800\text{ m}$ với vận tốc $60\text{ km/h}$. Thời gian đoàn tàu chạy qua hết cây cầu đó là bao nhiêu giây?
\shortans{60}
\loigiai{
Quãng đường đoàn tàu phải đi để qua hết cầu là:
$$800 + 200 = 1\,000\text{ (m)} = 1\text{ km}.$$
Vận tốc của tàu theo mét trên giây là:
$$60\,000 : 3\,600 = \dfrac{50}{3}\text{ (m/s)}.$$
Thời gian tàu chạy qua hết cầu là:
$$1\,000 : \dfrac{50}{3} = 60\text{ (giây)}.$$
}
\end{ex}

\begin{ex} %[5D10V61]
Hai thành phố cách nhau $180\text{ km}$. Hai xe ô tô khởi hành cùng một lúc từ hai thành phố và đi ngược chiều nhau. Sau $2$ giờ hai xe gặp nhau. Biết vận tốc xe đi từ A lớn hơn vận tốc xe đi từ B là $10\text{ km/h}$. Vận tốc của xe đi từ A là bao nhiêu ki-lô-mét trên giờ?
\shortans{50}
\loigiai{
Tổng vận tốc của hai xe là:
$$180 : 2 = 90\text{ (km/h)}.$$
Vận tốc của xe đi từ A là:
$$(90 + 10) : 2 = 50\text{ (km/h)}.$$
}
\end{ex}
